% !TEX program = pdflatex
\documentclass[11pt]{article}

% ---------- Encoding / Fonts ----------
\usepackage[T1]{fontenc}
\usepackage[utf8]{inputenc}
\usepackage{lmodern}
\usepackage{microtype}

% ---------- Layout ----------
\usepackage{geometry}
\geometry{margin=1in}

% ---------- Math / Theorems ----------
\usepackage{amsmath,amssymb,amsthm,mathtools}

% ---------- Hyperlinks ----------
\usepackage[hidelinks]{hyperref}

% ---------- Lists ----------
\usepackage{enumitem}

% ============================================================
% Theorem environments
% ============================================================
\theoremstyle{plain}
\newtheorem{theorem}{Theorem}
\newtheorem{proposition}{Proposition}
\newtheorem{lemma}{Lemma}
\newtheorem{corollary}{Corollary}

\theoremstyle{definition}
\newtheorem{definition}{Definition}
\newtheorem{assumption}{Assumption}

\theoremstyle{remark}
\newtheorem{remark}{Remark}

% ============================================================
% Title
% ============================================================
\title{\textbf{Topological Phases of Matter as Admissible Overlap Structures}\\
\large A projection--based account of robustness and bulk--boundary correspondence}
\author{Peter Nero}
\date{January, 2026}

\begin{document}
\maketitle

\begin{abstract}
Topological phases of matter exhibit robustness, quantized response, and protected
boundary modes that are not explained by local symmetry breaking or microscopic
Hamiltonian details. We present a structural account of topological phases based on
coarse--grained projection and global admissibility of local effective descriptions.
We show that phases correspond to admissible basins of coarse--grained dynamics,
and that topological phases arise precisely when locally well--defined effective
descriptions fail to admit a globally consistent completion. Topological invariants
appear as obstruction classes of overlap structures, explaining quantization and
robustness. Bulk--boundary correspondence follows inevitably: boundary modes arise
as compensating degrees of freedom required to restore local consistency at
interfaces. The framework applies equally to interacting and disordered systems
with a spectral gap or mobility gap and does not rely on band topology or
free--fermion assumptions.
\end{abstract}

\tableofcontents

% ============================================================
\section{Introduction}
% ============================================================

Topological phases of matter are distinguished by properties that are insensitive
to local perturbations, disorder, and microscopic details. These include quantized
transport coefficients, robust degeneracies, and protected edge or surface modes.
Unlike conventional phases, these features do not arise from symmetry breaking or
local order parameters.

Standard explanations invoke band topology, Berry curvature, or long--range
entanglement. While powerful, these descriptions are often model--specific and do
not clearly explain \emph{why} topological protection is inevitable or why bulk and
boundary phenomena are universally linked.

In this work we provide a structural explanation. We show that topological phases
arise because only certain global overlap structures of coarse--grained effective
descriptions are admissible. Topological invariants encode global consistency
constraints, and protected boundary modes appear when these constraints fail at
interfaces.

% ============================================================
\section{Minimal Structural Assumptions}
% ============================================================

We work entirely at the level of effective condensed--matter descriptions.

\begin{assumption}[Locality and stability]
\label{ass:locality}
The system admits a local description and possesses either a spectral gap or a
mobility gap ensuring stability under small local perturbations.
\end{assumption}

\begin{assumption}[Coarse--grained projection]
\label{ass:projection}
There exists a coarse--grained projection from microscopic degrees of freedom to
effective degrees of freedom capturing stable low--energy physics.
\end{assumption}

\begin{assumption}[Local effective descriptions]
\label{ass:local_eff}
On sufficiently small regions of space, the projected system admits a consistent
local effective description stable under perturbations.
\end{assumption}

No assumption is made about noninteracting particles, translational symmetry, or
specific lattice Hamiltonians.

% ============================================================
\section{Phases as Admissible Basins}
% ============================================================

\begin{definition}[Admissible basin]
An admissible basin is a region of the effective state space that is stable under
the projected dynamics and under small local perturbations.
\end{definition}

\begin{theorem}[Phase inevitability]
\label{thm:phases}
Any system satisfying Assumptions~\ref{ass:locality}--\ref{ass:local_eff} decomposes
into admissible basins corresponding to phases of matter.
\end{theorem}

\begin{proof}
Without such basins, arbitrarily small perturbations would destabilize macroscopic
properties, contradicting observed phase stability. Coarse--grained projection
selects stable regions of effective state space, which define phases.
\qedhere
\end{proof}

Conventional symmetry--breaking phases correspond to basins related by symmetry.
Topological phases correspond to basins distinguished by global constraints.

% ============================================================
\section{Topological Phases as Global Admissibility Constraints}
\label{sec:topo_admissibility}
% ============================================================

We now give a structural characterization of topological phases.

\subsection{Local effective descriptions and overlaps}

Let $\Sigma$ be the spatial manifold. Cover $\Sigma$ by overlapping regions
$\{U_\alpha\}$. On each $U_\alpha$, the system admits a local effective Hilbert
space $\mathcal{H}_\alpha$ and observable algebra $\mathcal{A}_\alpha$.

\begin{definition}[Overlap maps]
On overlaps $U_\alpha\cap U_\beta$, an overlap map
\[
\phi_{\alpha\beta} : \mathcal{H}_\alpha|_{U_\alpha\cap U_\beta}
\longrightarrow
\mathcal{H}_\beta|_{U_\alpha\cap U_\beta}
\]
identifies the two local descriptions. The overlap maps are required to preserve the local observable algebra and its
adjoint structure up to admissible equivalence.

\end{definition}

\subsection{Global admissibility}

\begin{definition}[Global admissibility]
\label{def:global_admissibility}
A collection $\{\mathcal{H}_\alpha,\phi_{\alpha\beta}\}$ is globally admissible if
the overlap maps satisfy the cocycle condition
\[
\phi_{\alpha\beta}\circ\phi_{\beta\gamma}\circ\phi_{\gamma\alpha}=\mathrm{Id}
\]
on all triple overlaps.
\end{definition}

\subsection{Topological obstruction}


\begin{theorem}[Topological phase criterion]
\label{thm:topo_phase}
A phase is topological if and only if, under the locality and gap assumptions
stated in Section~2, its locally admissible effective descriptions fail to admit
a globally admissible completion, with the failure classified by a nontrivial
topological obstruction.
\end{theorem}


\begin{proof}
If a globally admissible completion exists, the phase can be continuously
deformed to a trivial phase without closing the gap. If no such completion
exists, the failure is measured by a topological obstruction class that cannot
change under continuous local perturbations. \qedhere
\end{proof}

\begin{proposition}[Quantization of invariants]
Topological obstruction classes take values in discrete sets determined by the
topology of $\Sigma$ and the structure of overlap maps.
\end{proposition}

\begin{proof}
Obstruction classes live in integral cohomology or K--theory groups, which are
discrete. Continuous perturbations cannot change them without violating local
admissibility.
\qedhere
\end{proof}

% ============================================================
\section{Robustness Under Disorder and Perturbations}
% ============================================================

\begin{proposition}[Robustness of topological phases]
Topological phases are robust under all local perturbations that preserve
local admissibility and the gap or mobility gap.
\end{proposition}

\begin{proof}
Local perturbations deform overlap maps locally but cannot remove a global
obstruction without closing the gap or violating admissibility.
\qedhere
\end{proof}

This explains robustness to disorder and interactions.

% ============================================================
\section{Bulk--Boundary Correspondence}
\label{sec:bulk_boundary}
% ============================================================

Let $\Sigma_{\mathrm{bulk}}\subset\Sigma$ support a topological phase and let
$\partial\Sigma_{\mathrm{bulk}}$ denote its boundary with a trivial phase.

\begin{theorem}[Bulk--boundary correspondence]
\label{thm:bulk_boundary}
If the bulk phase carries a nontrivial global admissibility obstruction, then
there necessarily exist boundary--localized modes whose role is to restore
local overlap consistency at $\partial\Sigma_{\mathrm{bulk}}$.
These modes are protected against all perturbations preserving locality and
the bulk gap or mobility gap.
\end{theorem}

\begin{proof}
The bulk obstruction prevents global gluing of overlap maps. At the boundary,
local admissibility must still hold. This requires additional boundary degrees
of freedom to compensate for the obstruction. Because the obstruction is
topological, these modes cannot be removed without closing the gap.
\qedhere
\end{proof}

\begin{corollary}
If the bulk phase is globally admissible (topologically trivial), no protected
boundary modes are required.
\end{corollary}

\begin{remark}
This formulation applies equally to interacting, disordered, and superconducting
systems and does not rely on single--particle band topology.
\end{remark}

% ============================================================
\section{Relation to Known Topological Phases}
% ============================================================

The framework recovers known results:
\begin{itemize}
\item Chern insulators: obstruction is the first Chern class.
\item $\mathbb{Z}_2$ topological insulators: time--reversal--protected obstruction.
\item Topological superconductors: real--structure obstruction yields Majorana modes.
\item Interacting topological order: the obstruction no longer resides in
single-particle bundles but in the global consistency data of the emergent operator algebra, generalizing overlap maps to tensor-categorical fusion and braiding structures.
\end{itemize}

The unifying feature is global admissibility, not microscopic detail.

% ============================================================
\section{Scope and Limitations}
% ============================================================

The analysis applies to systems with a spectral gap or mobility gap. Gap closing
corresponds to basin boundary crossing and signals phase transitions. The framework
does not claim to classify all possible exotic phases, but explains the structural
origin of topological protection where it exists.

% ============================================================
\section{Conclusions}
% ============================================================

We have shown that topological phases of matter arise as admissible overlap
structures of coarse--grained effective descriptions. Topological invariants
measure global consistency obstructions, explaining quantization and robustness.
Bulk--boundary correspondence follows inevitably: protected boundary modes arise
to restore local consistency at interfaces.

This perspective unifies band topology, interacting topological order, disorder
robustness, and bulk--boundary correspondence within a single structural framework
and clarifies why topological phases are both rare and exceptionally stable.
\begin{thebibliography}{99}

\bibitem{MTT-Admissibility}
P.~Nero,
\emph{Modal Triplet Theory: Admissibility, Encodings, and the Structure of Physical Description},
Zenodo preprint, January 2026.
\url{https://doi.org/10.5281/zenodo.18255621}

\bibitem{MTT-Foundation}
P.~Nero,
\emph{Modal Triplet Theory: Foundation},
Zenodo preprint, September 2025.
\url{https://doi.org/10.5281/zenodo.16949762}

\bibitem{MTT-FixedPoints}
P.~Nero,
\emph{Fixed Points I--VI: Complete Coherence Spine},
Zenodo preprints, August 2025.
\url{https://doi.org/10.5281/zenodo.16948748}

\bibitem{MTT-ProjectionAdmissibility}
P.~Nero,
\emph{The Projection--Admissibility Principle: Structural Constraints on Effective Physical Description},
Zenodo preprint, January 2026.
\url{https://doi.org/10.5281/zenodo.18255838}

\bibitem{MTT-Closure}
P.~Nero,
\emph{Closure and Inevitability in Modal Triplet Theory},
Zenodo preprint, January 2026.
\url{https://doi.org/10.5281/zenodo.18255510}

\bibitem{MTT-CoherenceCapacity}
P.~Nero,
\emph{Coherence Capacity as the Fundamental Resource of Effective Physics},
Zenodo preprint, January 2026.
\url{https://doi.org/10.5281/zenodo.18255905}

\bibitem{MTT-CoherenceDynamics}
P.~Nero,
\emph{Dynamics of Coherence Capacity: Transport, Concentration, and Exhaustion},
Zenodo preprint, January 2026.
\url{https://doi.org/10.5281/zenodo.18256048}

\bibitem{MTT-QM}
P.~Nero,
\emph{Modal Triplet Theory: From MTT to Quantum Mechanics},
Zenodo preprint, September 2025.
\url{https://doi.org/10.5281/zenodo.17074246}

\bibitem{MTT-QFT}
P.~Nero,
\emph{From MTT to Quantum Field Theory},
Zenodo preprint, 2025.
\url{https://doi.org/10.5281/zenodo.17068816}

\bibitem{MTT-GR}
P.~Nero,
\emph{Modal Triplet Theory: From MTT to General Relativity},
Zenodo preprint, October 2025.
\url{https://doi.org/10.5281/zenodo.16950597}

\bibitem{MTT-UVQG}
P.~Nero,
\emph{Modal Triplet Theory: From MTT to a UV-Finite, Unitary Quantum Gravity},
Zenodo preprint, 2025.
\url{https://doi.org/10.5281/zenodo.17077671}

\bibitem{MTT-Measurement}
P.~Nero,
\emph{Measurement as Disturbance and Stabilization in Modal Triplet Theory},
Zenodo preprint, 2025.
\url{https://doi.org/10.5281/zenodo.17177404}

\bibitem{MTT-Irreversibility}
P.~Nero,
\emph{Projection, Probability, and Irreversibility: Shadow Bridges Between Measurement, Black Holes, and Cosmology in Modal Triplet Theory},
Zenodo preprint, January 2026.
\url{https://doi.org/10.5281/zenodo.18256408}

\bibitem{MTT-Bell-Beables}
P.~Nero,
\emph{Modal Fixed Points, Bell’s Beables, and the Limits of Factorization},
Zenodo preprint, 2025.
\url{https://doi.org/10.5281/zenodo.17076300}

\bibitem{MTT-Bell-Temporal}
P.~Nero,
\emph{Temporal Bell Inequalities and Global Consistency in Modal Triplet Theory},
Zenodo preprint, August 2025.
\url{https://doi.org/10.5281/zenodo.18208884}

\bibitem{MTT-Stochastic}
P.~Nero,
\emph{From Modal Triplet Theory to Indivisible Stochastic Processes: A First-Principles, Fully Rigorous Derivation},
Zenodo preprint, January 2026.
\url{https://doi.org/10.5281/zenodo.18254862}

\end{thebibliography}
\end{document}
